\setcounter{figcount}{1}

\section{Specify dataset}

This menu choice includes all of the procedures related to specifying
the dataset for the model.  A dataset must be specified {\em before}
any other step.

\subsection{Number of input nodes}

This menu choice allows you to set the number of input nodes.  Its
default is 1 input node, although that should {\em never} be used.

\subsection{Number of output nodes}

This menu choice allows you to set the number of output nodes.  Its
default is 1 output node.

\subsection{Load raw dataset}

This menu choice allows you to load a {\em raw} dataset from a file.
The file must be in the same directory as the program.  A raw dataset
is one whose inputs and outputs have not yet been normalized.
Generally, neural computational procedures require a dataset to be
normalized, whereas statistical procedures do not.  All datasets are
of the format:

\begin{center}
\begin{tabular}{cccccccc}
$i_{11}$ & $i_{12}$ & $\ldots$ & $i_{1I}$ & $o_{11}$ & $o_{12}$ &
$\ldots$ & $o_{1O}$ \\
$i_{21}$ & $i_{22}$ & $\ldots$ & $i_{2I}$ & $o_{21}$ & $o_{22}$ &
$\ldots$ & $o_{2O}$ \\
$i_{31}$ & $i_{32}$ & $\ldots$ & $i_{3I}$ & $o_{31}$ & $o_{32}$ &
$\ldots$ & $o_{3O}$ \\
$\ldots$ & $\ldots$ & $\ldots$ & $\ldots$ & $\ldots$ & $\ldots$ &
$\ldots$ & $\ldots$ \\
$i_{N1}$ & $i_{N2}$ & $\ldots$ & $i_{NI}$ & $o_{N1}$ & $o_{N2}$ &
$\ldots$ & $o_{NO}$ \\
\end{tabular}
\end{center}

Where $i$ is an input to the model, and $o$ is an output, and for $N$
examples (``exemplars'') in the set, with $I$ total inputs and $O$
total outputs.

These columns may be separated in the file by spaces, commas, or any
character which is non-numeric (1 through 9, `.', e, E, - or + are not
allowed as column separators.).

\subsection{Convert a raw dataset into a training set}

This menu choice will convert your raw dataset into a training set, by
normalizing its inputs and outputs to those specified in the menu
choice ``specify variate bounds and types.''  You will be asked if you
want to change the number of thresholds used to calculate ROC area.

You will also be asked if you want to save your training set into
a file.  {\bf This is generally recommended,} as you may want to load
the training set at some point in the future to resume training.
If you do save your training set, you will be asked for the filename
to save its scaling factors.  Scaling factors are used in the equation:
\begin{equation}\label{driverScale}
x_{norm} = lbound + S(x - min)
\end{equation}
where $x_{norm}$ is the normalized result, $lbound$ is the lower limit
of the normalized $x$ value (1 value for all), $S$ is the scaling
factor, and $min$ is the minimum value of each $x$ in the dataset. In
this way, should the model be used in the future, you will be able to
compute the normalized $x_{norm}$ value from the $x$ value that the
user supplies.

\subsection{Randomize a raw dataset into training and test sets}

This menu choice will randomly divide your raw dataset into training
and test sets.  As it will attempt to keep the fraction of outcomes
equal in both sets the same, the number of output nodes {\em must be}
1, and the output {\em must be} discrete (e.g. held to 0 or 1).  You
specify if an output is discrete in the menu choice ``specify variate
bounds and types.''

After the dataset has been randomized, you will be asked if you want
to change the number of thresholds used to calculate ROC area, and if
you want to save the newly created training and test sets to files.

As with converting a raw dataset into a training set, you will be
asked if you want to save your training set and scaling factors into a
file (see \hyperref[driverScale]{equation \ref{driverScale}}).  You
will also be asked if you want to save your test set into a file.
Again, {\bf this is generally recommended.}

You will be given the following choices to specify the number of
exemplars to be placed in the test set, with the remainder to be
placed in the training set (``$N1/N2$'' cross-validation):

\subsubsection{As a whole number}

This choice allows you to specify the number of exemplars in the test
set as a whole number.

\subsubsection{As a fraction with a numerator and denominator}

This choice allows you to specify the number of exemplars in the test
set as a fraction of the number of exemplars in the raw dataset.  You
will be asked for the numerator and the denominator of that fraction.

\subsubsection{As a decimal fraction}

This choice allows you to specify the number of exemplars in the test
set as a fraction of the number of exemplars in the raw dataset.  You
will be asked for the decimal notation of that fraction (e.g. $0.25$
would place $\frac{1}{4}$ of the exemplars from the raw dataset into
the test set, and the remainder into the training set.)

\subsection{Load training set}

This menu choice allows you to load a training set from a file.  The
file must be in the same directory as the program.  After the training
set has been loaded, you will be asked if you want to change the
number of thresholds used to calculate ROC area.

\subsection{Save training set}

This menu choice allows you to save your training set to a file, which
will be saved in the same directory as the program.  You will also be
asked for the filename to save its scaling factors.
See \hyperref[driverScale]{equation \ref{driverScale}} for an
explanation of scaling factors.

\subsection{Save scaling factors}

This menu choice allows you to save the scaling factors of the
training set, so that should the model be used in the future, you will
be able to compute the normalized $x_{norm}$ values from the $x$
values that the user supplies.  See \hyperref[driverScale]{equation
\ref{driverScale}} for an explanation of scaling factors.

\subsection{Load test set}

This menu choice allows you to load a test set from a file.  The file
must be in the same directory as the program.  After the test set has
been loaded, you will be asked if you want to change the number of
thresholds used to calculate ROC area.

\subsection{Save test set}

This menu choice allows you to save your test set to a file, which
will be saved in the same directory as the program.

\subsection{Log dataset operations to history file}

Set this option to ``yes'' if you want the operations you perform in
the ``Specify dataset'' menu to be logged to the history file {\bf
neuron.log}.  This file will appear in the same directory as the
program.

\subsection{Specify variate bounds and types}

In the following menu choices, variates are normalized according to
the equation:

\begin{equation}\label{USERinputScale}
x_{norm} = lbound + [({{x - min} \over {max - min}}) \times (hbound -
lbound)]
\end{equation}

Where $lbound$ is the lower boundary of a normalized value, $hbound$
is the upper boundary, $min$ is the minimum value of a variable {\em
in the training set}, $max$ is the maximum value of a variable {\em in
the training set}, $x$ is the variable's value and the normalized
value of $x$ is $x_{norm}$.

\subsubsection{Output is discrete}

Set this to ``yes'' if your outputs will always be either 0 or 1.

\subsubsection{Threshold for discrete output}

Set this to the number over which model outputs will be defined as 1,
and under which model outputs will be defined as 0 for discrete
output.  For example, if the threshold is $0.5$, a numeric model
output of $0.6$ will be transformed to 1, and a numeric model ouput of
$0.3$ will be transformed to 0.  The default is $0.5$.

\subsubsection{Lower limit of input variates}

Choose this menu item to set $lbound$ for inputs in equation
\ref{USERinputScale}.  The default is $-0.9$.

\subsubsection{Upper limit of input variates}

Choose this menu item to set $hbound$ for inputs in equation
\ref{USERinputScale}.  The default is $+0.9$.

\subsubsection{Lower limit of output variates}

Choose this menu item to set $lbound$ for outputs in equation
\ref{USERinputScale}.  The default is $+0.1$ if the output is not discrete
(and obviously 0 if it is).

\subsubsection{Upper limit of output variates}

Choose this menu item to set $hbound$ for outputs in equation
\ref{USERinputScale}.  The default is $+0.9$ if the output is not discrete
(and obviously 1 if it is).

\subsection{Number of thresholds to calculate ROC area}

Choose this menu item to set the number of thresholds over which ROC
area is calculated.  The output {\em must be} discrete, and the number
of outputs {\em must be} 1 for ROC area to be calculated.

\section{Specify model}

Choose this menu item to specify a {\em neural computational} model.
A dataset must be specified {\em before} this choice, and this menu
choice must obviously be specified before ``Use model''.  Discriminant
function analysis {\em does not} require this menu choice, and in
fact, it is useless for discriminant function analysis.

\subsection{Bias nodes}

Choose this menu item to specify whether or not bias nodes will be
included in your neural computational model.  In general, for a
real-world model, you should {\em always} include bias nodes.  The
default is ``yes''.

\subsection{Number of hidden layers / nodes on each layer}

Choose this option to specify the number of hidden layers and number
of nodes on each layer.  The default is 1 layer with 1 node, although
neural computational models will generally have $>$ 1 hidden node.
Most models will, however, only have 1 hidden layer.

\subsection{Output error function}

Choose this option to specify the output error function.  ``LMS'' is
least-mean-squared:

\begin{equation}
E = \frac{1}{2} \sum_{i=1}^{m} (y_i - o_i)^2
\end{equation}
where $E$ is the error, $m$ the number of outputs, $y$ is the known
value and $o$ the model's output.

You may also choose the cross-entropy error function ``X-entropy'':

\begin{equation}
E = \sum_{i=1}^{m} -y\log(o)-(1-y)\log(1-o)
\end{equation}

For the cross-entropy error function, the output {\em must be}
discrete.  Note that for sigmoidal output nodes, where for
output $\vec{\bf{o}}$,

\begin{displaymath}
\vec{\bf{o}} = \frac{1}{1+e^{-\vec{\bf{x}}}}
\end{displaymath}

if any element of $\vec{\bf{x}}$ is very large or very small,
the program will round to $0$ or $1$, and the $\log$ terms in the
cross-entropy error function will cause an error.  In such cases,
an approximation is made (see the document {\bf manifest.pdf}
for a detailed explanation of that approximation), and the following
error will be generated:

\begin{verbatim}
WARNING: Numerical out of bounds encountered when calculating error
\end{verbatim}

\subsection{Load network from a file}

This menu choice allows you to load a neural computational model from
a file.  The file must be in the same directory as the program.

\subsection{Binary logistic regression}

This menu choice specifies binary logistic regression, which is
mathematically equivalent to a single output node neural network with
a bias, the sigmoidal transfer function, no hidden nodes, and the
cross-entropy error function. {\em \beginred Always run a binary
logistic regression model} {\bf before} {\em starting to build a
neural network. \endcolor}  If your binary logistic model has
acceptable goodness-of-fit, you don't need to build a neural network.

Note that at the end of a binary logistic run, $\beta$ weights will be
reported, along with their standard errors and $p$-values for the Wald
test that the $\beta$ weights are nonzero.\footnote{see Hosmer and
Lemeshow, Applied Logistic Regression, 2nd edition, John Wiley \&
Sons, Inc., New York, 2000, pp 36--42}  Note also that the indexing
begins with input number 0, which is often referred to as the $y$-
intercept.  {\em \beginred In this case, $\beta_0$ is} {\bf not} {\em
the $y$-intercept. \endcolor}  It is the $\beta$ weight of the first
variate, indexed to 0.  The $y$-intercept is clearly labeled as
``$y$-intercept'' in the report table.

\subsection{Log last operation to file}

Set this option to ``yes'' (the default) if you want the last neural
computational operation performed to be logged to the file {\bf
model.txt}.  This file will be {\em overwritten} with each run, and
will appear in the same directory as the program.

\subsection{Log to history file}

Set this option to ``yes'' (the default) if you want the neural
computational operations that are performed logged to the file {\bf
neuron.log}.  This file will be {\em appended} with each operation,
and will appear in the same directory as the program.

\section{Use model}

Choose this menu item to use a specified {\em neural computational}
model.  A dataset must be specified {\em before} this choice, and the
``Specify Model'' menu must also have been entered before ``Use
model''.

\subsection{Randomize weights}

Choose this option to randomize the weights of the neural
computational model.  Prior to training the model, weights {\em must
be} set, either by randomization or by loading the model from a file
in the ``Specify model'' menu.  By default, weights are randomized
between $-0.5$ and $+0.5$.

\subsection{Set learning rate}

Choose this option to set the learning rate $\eta$ of the neural
computational model.  The default is $\eta = 1/N$ for on-line
backpropagation (where $N$ is the number of exemplars), and $\eta = 1$
for off-line (``batch'' or ``epoch'') training.

You may choose the learning rate to be ``automatic'' (see the Manifest
documentation for a fuller explanation of automatic stepsize
selection.) Briefly, automatic stepsize selection is designed to
prevent oscillation around a local error minimum.  Due to the increase
in time required for this feature, it may be most optimal to train to
near a local error minimum before setting automatic stepsize selection
{\em on}.

\subsection{Set stopping conditions}

Choose this option to set conditions by which the neural computational
model will stop training.  These include:

\subsubsection{Set maximum number of iterations}

The model will stop training after the specified number of iterations.
The default is 1000000, {\em which is arbitrary.  Do not consider any
neural computational model to be trained by stopping at any arbitrary
threshold.  You} must {\em inspect the model yourself to see if it has
stopped training.}

\subsubsection{Lower error limit}

Use this choice to set the lower error limit that, if the error of the
model decreases below, will result in cessation of training.  For
example, you could set this to the limit of your compiler.  By
default, no lower error limit is set.

\subsubsection{Limit of change in error over 1 iteration}

Use this choice to set the lower error limit that, if the change in
error over 1 iteration is below, will result in cessation of training.
Thus, in

\begin{equation}
\Delta E = E_{t+1} - E_t
\end{equation}

Where $E_t$ is the error after a training iteration, and $E_{t+1}$ is
the error after the next training iteration, if $\Delta E$ is less
than the specified value, training stops.  By default, this condition
is not specified.

\subsubsection{Stop if increase over error window width}

Allows you to set a number of iterations that after which the error
increases, training will stop.  Thus, in

\begin{equation}
E_{n+\Delta} > E_{n}
\end{equation}

Where $n$ is the current iteration, $E$ the error at an iteration, and
$\Delta$ the error window width, training will stop if the inequality
is satisfied.  By default, this condition is not specified.

\subsubsection{Maximum absolute gradient limit}

All neural networks update their weights in the following form:

\begin{equation}
{\bf w}_{t+1} = {\bf w}_t - \eta {\bf g}_t
\end{equation}

where $\bf{w}_t$ is the weight structure at time $t$, $\eta$ is the
learning rate, and $\bf{g}_t$ is the gradient.  This menu choice
allows you to stop training when the maximum absolute value of the
elements in the gradient structure $\bf{g}_t$ decreases below the
chosen limit.  The default is $10^{-6}$.

\vspace{9 pt}{\bf \large \beginred READ THIS! \endcolor}

\beginred The maximum absolute gradient limit should be the primary
means by which you judge that training has ended, {\em not} by
the error minimum. \endcolor

For canonical backpropagation, using the maximum absolute gradient
limit will significantly slow training (not by the number of
iterations it takes to train, but by the amount of computational time
consumed per iteration.)  Thus, one strategy is to turn the maximum
gradient choice off until you are near a local error minimum as judged
loosely by the decrease in the error, then to turn on this menu
choice.

Don't be surprised if the maximum absolute gradient fluctuates
(increases and decreases.)  Remember, there is more than one node and
one layer in even the simplest neural network.  Keep an eye on the
overall trend of the maximum absolute gradient to insure it is
decreasing.

\subsection{Batch Epoch}

Choose this option to set batch/epoch learning.

The {\em on-line} backpropagation learning algorithm is given by the
formula:
\begin{equation}
{\bf w}_{t+1} = {\bf w}_t - \eta {\bf g}_k({\bf w}_t)
\end{equation}
where $k$ is the exemplar index, ${\bf w}$ is the weight, $t$ the
iteration, ${\bf g}$ is the error gradient with respect to the weight,
and $\eta$ is the learning rate.  By default, if on-line learning is
chosen, $\eta = 1/N$, where $N$ is the number of exemplars.

The {\em off-line}, also known as ``batch'' or ``epoch''
backpropagation learning algorithm is given by the formula:
\begin{equation}
{\bf w}_{t+1} = {\bf w}_t - \eta {\bf g}({\bf w}_t)
\end{equation}
where the {\em average gradient} ${\bf g}$ is defined by the formula:
\begin{equation}
{\bf g} = (1/N) \sum_{k=1}^N {\bf g}_k.
\end{equation}
By default, if off-line learning is chosen, $\eta = 1$.

See the neUROn2++ Design Manifest for a more detailed description of
on-line and off- line learning.

If batch/epoch is set {\em on}, off-line, ``batch'' or ``epoch''
learning is used. If batch/epoch is set {\em off}, on-line learning is
used.

\subsection{Weight Decay}

The {\em error function}, $E_k$, for the $k$th exemplar is
defined by the formula:
\begin{equation}\label{USERdeferrorfunk}
E_k({\bf w}) =
e(y^k, {\bf o}^{(m)}_k) +
\frac{1}{2\sigma_w^2} \sum_{i=1}^m |{\bf w}|^2.
\end{equation}

where $y$ is the known or target outcome, and ${\bf o}$ is the model's
output.  The second term on the right hand side of
(\ref{USERdeferrorfunk}) is required to guarantee reliable statistical
inference. It basically `` stabilizes'' the learning process by
preventing the weights in the network from becoming too large or too
small.  $\sigma_w$ is the standard deviation of the weights {\em \bf
chosen by the user.}  Thus, a connection weight is expected to lie
between $-1.96 \sigma_w$ and $+1.96 \sigma_w$ with probability equal
to about 95\%. Adjustment of this parameter can ``change results'' of
the statistical analysis so $\sigma_w$ is constrained to be exponents
of 10 (e.g. 10, 100, 1000, $\ldots$ $10^{12}$.  Once chosen,
$\sigma_w$ should not be changed.

A reasonable choice for $\sigma_w$ might be $\sigma_w = 100$. This
basically constrains the ``learning process'' to find connection
strengths for the ANN such that the connection strength weight ranges
between $-200$ and $+200$. The choice of $\sigma_w$ should be reported
in conjunction with all statistical analysis results. Note that as
$\sigma_w \rightarrow \infty$, the effect of the summation term
becomes negligible.

The learning rate equivalent of equation \ref{USERdeferrorfunk} is

% Note the necessity of \left. even though it doesn't show
\begin{eqnarray*}
\vec w_{t+1} & = & \vec w_t - \eta \left. \D{E'}w \right|_{w_t} \\
& = & \vec w_t - \eta \left. \D{E}w \right|_{w_t} - 2\eta\lambda\vec w_t \\
& = & (1-2\eta\lambda)\vec w_t - \eta \left. \D{E}w \right|_{w_t}
\end{eqnarray*}

where $(1-2\eta\lambda)$ is ``weight decay''. $\lambda$ may
be set by either of 2 methods:

\begin{itemize}

\item The complexity of the population objective function $>>$ sample
objective function:

\begin{equation}
\lambda = \frac 1{N2\sigma_w^2}
\end{equation}

\item The complexity of the population objective function $\simeq$
sample objective function:

\begin{equation}\label{USERsameComplexEqn}
\lambda = \frac 1{2\sigma_w^2}
\end{equation}

\end{itemize}

In general, as we don't know the complexity of the population
objective function a priori, we choose equation \ref{USERsameComplexEqn}
to calculate $\lambda$.

The following choices allow you to specify weight decay:

\subsubsection{Don't change anything}

Allows the uninformed user (which, assuming you've read this section,
you are {\em not},) to exit without changing the way in which weight
decay is calculated.

\subsubsection{lambda = 1 / ( 2 * sigma$^\wedge$2 )}

Specifies that the weight decay term be calculated by

\begin{displaymath}
\lambda = \frac 1{2\sigma_w^2}
\end{displaymath}

and sets $\sigma_w$ to be an exponent of 10 from $10^1$ to $10^{12}$.

\subsubsection{lambda = 1 / ( 2 * N * sigma$^\wedge$2 )}

Specifies that the weight decay term be calculated by

\begin{displaymath}
\lambda = \frac 1{N2\sigma_w^2}
\end{displaymath}

and sets $\sigma_w$ to be an exponent of 10 from $10^1$ to $10^{12}$.

\subsubsection{No weight decay}

Turns off weight decay.

\vspace{9 pt}{\bf \large \beginred READ THIS! \endcolor}

Using too large of a weight decay term will degrade network modeling
performance, but it is critical to use weight decay in order for
statistical inference, for example, stepwise regression.  Thus,
setting a proper weight decay value is of utmost importance.  Here is
a good strategy to determine the proper weight decay value:

{\em (Note that the default value for the weight decay term is
$\frac 1{2\sigma_w^2}$ with $\sigma_w = 100$.  This is to prevent the
uninformed user from making incorrect statistical inferences.
However, you are no longer the uninformed user. In the first step, you
will turn off weight decay.)}

\begin{enumerate}

\item Turn weight decay off.

\item Experiment with different neural network architectures until you
find one that is as near as possible a strict local minimum, and with
as high ROC area as possible.

\item Save the neural network from Step 2.

\item Load the neural network from Step 2, and turn on weight decay.
Use the choice for $\lambda = \frac 1{2\sigma_w^2}$, and choose a low
initial value for $\sigma$, for example $10^2$.

\item Train the neural network to as near a possible a strict local
minimum (you'll be starting with the final weights from Step 2.)

\item Compare the ROC area to that from Step 2.  If it is similar, go
to step 8.

\item If the ROC area has substantially degraded, increase $\sigma$,
load the neural network from Step 2, and go to Step 5.

\item Save the neural network, and record $\sigma_w$, as you will cite
it when you publish your results. Perform stepwise regression with
this value in the weight decay term, beginning with the final weights
from the network you saved in this Step.

\end{enumerate}

\subsection{Print counter}

This menu choice allows you to indicate whether printing will be
linear (prints after a set number of iterations) or logarithmic.
Logarithmic is recommended.

\subsection{Train model}

This menu choice initiates neural computational model training.  You
will choose the type of gradient descent.  For an explanation of
conjugate gradient descent and the Shanno algorithm, see Golden, pp.
217--218 and 221--222.\footnote{Richard M. Golden, Mathematical
Methods for Neural Network Analysis and Design, 1996, MIT Press, ISBN
0-262-07174-6, pp. 217--218, 221--222}  A brief explanation of these
methods may also be found in the design manifest.

The menu choices are ordered by robustness of the algorithm, from less
to more.  Thus, conjugate gradient descent is less robust than the
Shanno algorithm.

\vspace{9 pt}{\bf \large \beginred READ THIS! \endcolor}

Conjugate gradient descent and the Shanno algorithm generally are much
faster at finding a local error minimum than canonical backpropagation
{\em once one is near the local error minimum}.  Thus, your training
strategy should be to use canonical backpropagation until you are near
a local error minimum (as judged by the maximum absolute gradient),
and then apply either conjugated gradient descent or the Shanno
algorithm. Empirically try each advanced algorithm, and choose the one
to use based on your trials.

\subsubsection{Canonical backpropagation}

This menu choice implements standard, plain vanilla backpropagation.

\subsubsection{Conjugate gradient descent}

This menu choice implements conjugate gradient descent.  See the
design manifest and Golden pp. 221--222.

\subsubsection{Shanno algorithm}

This menu choice implements the Shanno algorithm.  See the design
manifest and Golden pp. 217--218.

\subsection{Save the network to a file}

This menu choice allows you to save your trained neural computational
model with its weights to a file.  You may train the model repeatedly
prior to saving it, if you so choose.  If you have one output, and it
is discrete, you will also be asked if you want to save the network's
guesses to a file. These guesses will be saved as a 2 column file,
with the first column being the known value, and the second column
being the network's guess. You will be asked if you want to save
guesses for the training set, and if it is loaded, the test set.

\subsection{Stepwise regression}\label{stepregress}

Stepwise regression allows you to determine which input variables are
statistically significant to your trained neural computational model.
For a more complete discussion, see the neUROn2++ Design Manifest.
If you have not just finished training a neural network, you will be
asked for the final error of the network model to regress.

\subsubsection{Specify variables}

A variable may be represented by more than 1 input node.  For example,
ethnicity may be represented by 4 nodes: African-American, Asian,
Hispanic and Caucasian.  Or a variable may be represented by 2 nodes:
1 in which the value is present or absent, and the other by the value
of the variable itself.  Variable structure may be input from the
command line or loaded from file.

Input characters include:
\begin{itemize}
\item {\bf ;} delimits variables
\item {\bf ,} delimits nodes
\item {\bf -} specifies a range of nodes
\item spaces are ignored
\end{itemize}

For example, ``0-3,5; 4,8; 6,7'' indicates that variable 0 is composed
of input nodes 0,1,2,3 and 5, variable 1 is composed of input nodes 4
and 8, and variable 2 is composed of input nodes 6 and 7.

\subsubsection{Stepwise reverse regression}

Stepwise reverse regression proceeds by training the full network, and
recording that network's final error.  Input variables are removed 1
at a time, the subnetworks so generated are trained to completion and
their errors recorded, and the variable with the largest $p$-value
(least significant) is then recorded with its $p$-value.  Using the
final error of the network which is deficient the one node with the
largest $p$-value as a starting point, input variables are then
removed 2 at a time, with each pair containing the 1 variable with the
largest $p$-value and 1 of the remaining input variables.  The pair of
variables with the largest $p$-value (least significant) is then
recorded with its $p$-value, and the process repeats until a chosen
threshold $p$-value is reached.

Figure~\arabic{figcount} demonstrates stepwise reverse regression.

\begin{figure}[htb]\addtocounter{figcount}{1}
\begin{picture}(4.2,2.2)
\put(0,0){\framebox(4.2,2.2)}
% for dvi
\put(0.1,0.1){\includegraphics[width=4in,height=2in]{reverse.bmp}}
% for dvipdfm
\put(0.1,0.1){\special{pdf: epdf width 4in height 2in (reverse.pdf)}}
\end{picture}
\makebox[4in]{\rule{0in}{2in}}
\caption{Stepwise reverse regression}
\end{figure}

\subsubsection{Stepwise forward regression}

Stepwise forward regression proceeds by training a baseline network
with no input nodes, and recording that network's final error.  Input
variables are added 1 at a time, the subnetworks so generated are
trained to completion and their errors recorded, and the variable with
the smallest $p$-value (most significant) is then recorded with its
$p$-value.  Using the final error of the network which is the one node
with the smallest $p$-value as a starting point, input variables are
then added 2 at a time, with each pair containing the 1 variable with
the smallest $p$-value and 1 of the remaining input variables.  The
pair of variables with the smallest $p$-value (most significant) is
then recorded with its $p$-value, and the process repeats until a
chosen threshold $p$-value is reached.

Figure~\arabic{figcount} demonstrates stepwise forward regression.

\begin{figure}[htb]\addtocounter{figcount}{1}
\begin{picture}(4.2,2.2)
\put(0,0){\framebox(4.2,2.2)}
% for dvi
\put(0.1,0.1){\includegraphics[width=4in,height=2.1in]{forward.bmp}}
% for dvipdfm
\put(0.1,0.1){\special{pdf: epdf width 4in height 2.1in(forward.pdf)}}
\end{picture}
\makebox[4in]{\rule{0in}{2.1in}}
\caption{Stepwise forward regression}
\end{figure}

\section{Discriminant function analysis}

Choose this menu item to use discriminant function analysis to model
your dataset(s).

\subsection{Linear discriminant function analysis }

Choose this menu item to model your dataset(s) with {\em LDFA}, linear
discriminant function analysis.

\subsection{Quadratic discriminant function analysis }

Choose this menu item to model your dataset(s) with {\em QDFA},
quadratic discriminant function analysis.

\subsection{Log last operation to file}

Set this option to ``yes'' (the default) if you want the last
statistical operation performed to be logged to the file {\bf
model.txt}.  This file will be {\em overwritten} with each modeling
procedure, and will appear in the same directory as the program.

\subsection{Log to history file}

Set this option to ``yes'' (the default) if you want the statistical
operations that are performed logged to the file {\bf neuron.log}.
This file will be {\em appended} with each operation, and will appear
in the same directory as the program.
